%%% Preamble (styling follows docs/manual/main.tex)
\documentclass[paper=a4, fontsize=11pt]{scrartcl}
\usepackage[T1]{fontenc}
\usepackage{fourier}

\usepackage[english]{babel}
\usepackage[protrusion=true,expansion=true]{microtype}
\usepackage{amsmath,amsfonts}
\usepackage{url}
\usepackage{hyperref}
\usepackage{booktabs}
\usepackage{minted}  % compile with -shell-escape

\usepackage{sectsty}
\allsectionsfont{\centering \normalfont\scshape}

%%% Page numbers only
\usepackage{fancyhdr}
\pagestyle{fancyplain}
\fancyhead{}
\fancyfoot[L]{}
\fancyfoot[C]{}
\fancyfoot[R]{\thepage}
\renewcommand{\headrulewidth}{0pt}
\renewcommand{\footrulewidth}{0pt}
\setlength{\headheight}{13.6pt}

%%% Equation and float numbering
\numberwithin{equation}{section}
\numberwithin{table}{section}

%%% Formatting parameters
\setlength{\parindent}{0pt}
\setlength{\parskip}{6pt}
\setlength{\voffset}{-0.5cm}
\setlength{\hoffset}{-0.5cm}
\setlength{\textheight}{24.5cm}
\setlength{\textwidth}{16cm}
\emergencystretch=4em

%%% Colours and typewriter names, as in the manual
\usepackage[dvipsnames]{xcolor}
\newcommand{\codecolor}{\color{BrickRed}}
\newcommand{\filecolor}{\color{NavyBlue}}
\usepackage{xstring}
\newcommand{\brk}[1]{\leavevmode\begingroup\ttfamily
  \def\_{\textunderscore\allowbreak}%
  \StrSubstitute{#1}{/}{/\allowbreak}[\brkA]%
  \StrSubstitute{\brkA}{.}{.\allowbreak}[\brkA]%
  \brkA\endgroup}
\newcommand{\File}[1]{\leavevmode{\filecolor\brk{#1}}}
\newcommand{\Code}[1]{\leavevmode{\codecolor\brk{#1}}}
\definecolor{codegray}{rgb}{0.5,0.5,0.5}
\definecolor{backcolour}{RGB}{253,246,228}
\definecolor{linkteal}{RGB}{0,128,128}
\hypersetup{colorlinks=true, linkcolor=linkteal, citecolor=linkteal, urlcolor=linkteal}

\usemintedstyle{abap}
\setminted{bgcolor=backcolour,
           baselinestretch=1.1,
           linenos,
           numbersep=4pt}
\renewcommand{\theFancyVerbLine}{\color{codegray}{\tiny\arabic{FancyVerbLine}}}
\newcommand{\fort}[1]{\smash{\mintinline{Fortran}{#1}}}

\newcommand{\eps}{\varepsilon}

%%% Title
\newcommand{\horrule}[1]{\rule{\linewidth}{#1}}
\title{
  \normalfont \normalsize \textsc{OpenMULTI3D technical note} \\ [25pt]
  \horrule{0.5pt} \\[0.4cm]
  \Large Single-precision formal solvers \\
  \horrule{2pt} \\[0.5cm]
}
\author{\normalfont \normalsize October 2026}
\date{}

\begin{document}
\maketitle

\begin{abstract}
The linear, quadratic and Hermite (cubic B\'ezier) long-characteristic formal
solvers in \File{SRC/rt\_solvers\_mod.f90} evaluate their interval weights in
single precision (fp32) without data-dependent branches. Every weight is
accurate to a few units in the last place (ulp) over the whole range of optical
depths that occurs, including the bounded maser intervals with negative optical
depth. This note gives the motivation, the formulation and its accuracy.
\end{abstract}

\section{Motivation}

Across one interval of a characteristic, the formal solution is
\begin{equation}
  I_c = I_u\, e^{-\Delta} + \int_0^{\Delta} S(t)\, e^{-t}\, \mathrm{d}t ,
  \label{eq:formal}
\end{equation}
where $\Delta$ is the optical depth of the interval, $t$ is measured backward
from the target point $c$, and $u$ denotes the upwind point. Each scheme
replaces $S(t)$ by an interpolant in a few source values and slopes, so that
$I_c$ is a weighted sum of these with weights that depend only on $\Delta$.
The weights are evaluated once per frequency and grid point, which makes them
the innermost work of every solver.

The opacities, source functions and intensities are stored in fp32. The
weights were nevertheless formed in fp64, because their textbook expressions
cancel: $1-e^{-\Delta}$ loses all relative accuracy as $\Delta\to0$, and the
moment combinations behind the higher-order weights cancel to
$O(\Delta^{n})$ with $n$ up to 4. The fp64 code therefore also switched to a
truncated Taylor series below a threshold and clamped the exponential at
$\Delta=80$. Both the branches and the fp64 arithmetic halve or remove the SIMD
width of the frequency loop, and in fp64 the arithmetic cost is no longer
negligible against the memory traffic of the fp32 arrays. The previous linear
step, for example, read

\begin{minted}{Fortran}
dt = dsup05*(xxup + xx)
dt = max(dt, -maser_dtau_max)
if (dt < 80) then
  ex = exp(-dt)
  if (abs(dt) > 1d-3) then
    u0 = 1 - ex
    Psi = 1 - u0/dt
  else
    Psi = dt*(0.5 - 0.16666667*dt)
    u0 = dt*(1 - Psi)
  endif
  Inu = Inuup*ex + (u0 - Psi)*ssup + Psi*ss
else
  u0 = 1/dt
  Psi = 1 - u0
  Inu = u0*ssup + Psi*ss
endif
\end{minted}

The approach below removes every cancellation analytically, so that fp32
suffices. All forms are written so that both sides of a switch are
evaluated and one is selected without a branch.

\section{The shared exponential}

All three schemes need $E=e^{-\Delta}$ and $U_0=1-e^{-\Delta}$, each to full
relative accuracy, for $\Delta$ between the maser cap ($-2$ by default) and
the optically thick limit. They are computed together by range reduction
(\fort{lw_exp}). With
\begin{equation}
  \Delta = n\ln 2 + r,\qquad n=\operatorname{nint}(\Delta/\ln2),\qquad |r|\le \tfrac12\ln2 ,
\end{equation}
where $\ln2$ is split into an exactly representable leading part and a small
correction (Cody--Waite), so that $r$ is exact to rounding, and
\begin{equation}
  s = 2^{-n},\qquad m = 1-e^{-r} = r\,q(r),
\end{equation}
one has
\begin{equation}
  E = s\,(1-m),\qquad U_0 = (1-s) + s\,m .
  \label{eq:EU0}
\end{equation}
Here $s$ is exact (it is written directly into the exponent field), $q$ is a
relative-error minimax polynomial of degree 5 on $|r|\le\frac12\ln2$, and
neither $E$ nor $U_0$ is obtained by subtracting the other from 1. For
$\Delta\to0$, $n=0$ and $U_0 = m$ carries full relative accuracy; for large
$\Delta$, $1-s$ is exact. $\Delta$ is clamped at 87 inside the exponential
only, where $E$ falls below every quantity it multiplies.

In \fort{lw_exp}, with \fort{q(r)} the degree-5 polynomial evaluated by Horner's
rule,
\begin{minted}{Fortran}
tc = min(t, 87.0)
n  = anint(tc*LOG2E)
r  = (tc - n*C1) - n*C2              ! C1 + C2 = ln 2, n*C1 exact
m  = r*q(r)                          ! 1 - exp(-r)
s  = transfer(ishft(127 - int(n), 23), 1.0)   ! 2**(-n) from its exponent bits
E  = s - s*m
u0 = (1.0 - s) + s*m
\end{minted}

\section{Branch-free switching}

Each weight has a form that is accurate for small $|\Delta|$ (a series) and one
that is accurate above some threshold (closed form in $E$ and $U_0$). Both are
evaluated at every point, and one is selected,
\begin{equation}
  w = \chi\, w_{\mathrm{series}} + (1-\chi)\, w_{\mathrm{closed}},\qquad
  \chi = [\,\text{series range}\,],
\end{equation}
implemented by \fort{lw_select(a, b, first)} as a bitwise mask on the two fp32
values,
\begin{minted}{Fortran}
sel = merge(-1, 0, first)
lw_select = transfer(ior(iand(transfer(a, 1), sel),      &
                         iand(transfer(b, 1), not(sel))), 1.0)
\end{minted}
so that a non-finite value in the unused form never enters the result. The closed
forms are evaluated at an argument clamped into their own range where that
keeps them finite.

The series are relative-error minimax polynomials, fitted on the stated
interval by iteratively reweighted least squares, rounded to fp32 and then
verified in compiled code against a quad-precision reference.

\section{Linear scheme}

With $S$ linear in $t$ between $S_c$ ($t=0$) and $S_u$ ($t=\Delta$),
\begin{equation}
  I_c = E\, I_u + W\, S_u + \Psi\, S_c ,\qquad
  \Psi = 1-\frac{U_0}{\Delta},\qquad W = U_0 - \Psi = \frac{U_0}{\Delta}-E .
  \label{eq:linear}
\end{equation}
$\Psi$ is also the diagonal of the approximate $\Lambda$ operator (divided by
$\chi_c$), and the preconditioned rates need the escape factor $1-\Psi$.

$\Psi$ cancels as $\Delta\to0$ and $1-\Psi$ cancels as $\Delta\to\infty$. Hence
\begin{equation}
  \Psi = \begin{cases}
    \Delta\, P(\Delta), & |\Delta|\le 1,\\[2pt]
    1-U_0/\Delta, & \text{otherwise},
  \end{cases}
  \qquad
  P(\Delta) = \sum_{k\ge0}\frac{(-\Delta)^k}{(k+2)!},
\end{equation}
with $P$ a minimax polynomial of degree 7 on $|\Delta|\le1$. Above $|\Delta|=1$
the escape factor is returned as $U_0/\Delta$ and $W$ as $U_0/\Delta-E$, so that
neither the $W\sim1/\Delta$ tail nor $1-\Psi$ is formed by subtraction from 1.
Every output is within 4 fp32 ulp of the exact value for
$-2\le\Delta\le90$.

In \fort{linear_weights}, with \fort{lw_series(t)} $=\Delta P(\Delta)$,
\begin{minted}{Fortran}
call lw_exp(t, E, u0)
ut = u0/t
ps = lw_series(t)
thin = abs(t) <= 1.0
Psi = lw_select(ps, 1.0 - ut, thin)
Esc = lw_select(1.0 - ps, ut, thin)  ! escape factor 1 - Psi
W   = lw_select(u0 - ps, ut - E, thin)
\end{minted}
and the step in the frequency loop is
\begin{minted}{Fortran}
dt = dsup05*(xxup + xx)
dt = max(dt, -real(maser_dtau_max, 4))
call linear_weights(dt, E, W, Psi, Esc)
Inu = Inuup*E + W*ssup + Psi*ss
\end{minted}

The step itself is then $I_c = E\,I_u + W S_u + \Psi S_c$, followed by the maser
caps on $I_c$ and the selection of $\Psi=0$ (no preconditioning) where
$\Delta\le0$ or $\chi_c\le0$, all without branches.

\section{Quadratic scheme}

The quadratic solver is the parabolic (B\'ezier-equivalent) scheme of Multi3D.
Measure optical depth from the upwind point, $s=\Delta-t$, and represent the
source by the parabola through $S_u$ and $S_c$ with slope $S'_c=\mathrm{d}S/\mathrm{d}\tau$ at
the target (MC-limited from the one-sided slopes on both sides of $c$). With
the moments
\begin{equation}
\begin{gathered}
  U_0 = \int_0^\Delta e^{-(\Delta-s)}\mathrm{d}s = 1-E,\qquad
  U_1 = \int_0^\Delta s\,e^{-(\Delta-s)}\mathrm{d}s = \Delta-U_0,\\
  U_2 = \int_0^\Delta s^2 e^{-(\Delta-s)}\mathrm{d}s = \Delta^2-2U_1,
\end{gathered}
\end{equation}
and $q=U_2/\Delta$, Eq.~(\ref{eq:formal}) becomes
\begin{equation}
  I_c = E\,I_u + A\,S_u + B\,S_c + C\,S'_c,\qquad
  B = \frac{2U_1-q}{\Delta},\quad A = U_0-B,\quad C = q-U_1 .
\end{equation}
On the last plane of a sweep the stencil is shifted one cell upwind and the
slope sits at the upwind end; the same moments then give
$A=q/\Delta$, $B=U_0-A$, $C=U_1-q$.

$U_1\simeq\Delta^2/2$ and $q\simeq\Delta^2/3$ cancel as $\Delta\to0$.
$U_1 = \Delta\Psi$ with the $\Psi$ of the linear scheme, so it inherits that
treatment. For $q$,
\begin{equation}
  q = \begin{cases}
    \Delta^2 R(\Delta), & |\Delta|\le2,\\[2pt]
    \Delta - 2U_1/\Delta, & \text{otherwise},
  \end{cases}
  \qquad
  R(\Delta) = \sum_{k\ge0}\frac{2(-\Delta)^k}{(k+3)!},
\end{equation}
with $R$ a minimax polynomial of degree 8 on $|\Delta|\le2$. The series range
extends to 2 rather than 1 because the closed form still loses about 10 ulp near
$\Delta=-1.5$. $E$, $U_0$, $U_1$ and $q$ are all within 4 ulp of the exact values
for $-2\le\Delta\le87$ (\fort{parab_moments}).

In \fort{parab_moments}, with \fort{pm_r(t)} $=R(\Delta)$,
\begin{minted}{Fortran}
call lw_exp(t, E, U0)
U1 = lw_select(t*lw_series(t), t - U0, abs(t) <= 1.0)
q  = lw_select(t*t*pm_r(t), t - 2.0*(U1/t), abs(t) <= 2.0)
\end{minted}
and in the frequency loop of \fort{long_char_quadratic}
\begin{minted}{Fortran}
rdt = 1.0/dt
call parab_moments(dt, ex, U0, U1, q)
B = (2.0*U1 - q)*rdt                 ! on ssc, the target
A = U0 - B                           ! on ssup, the upwind point
C = q - U1                           ! on dsdtau
Inu = Inuup*ex + A*ssup + B*ssc + C*dsdtau
\end{minted}

The weights $A$, $B$, $C$ are then formed in fp32. At large $\Delta$,
$A=U_0-B\simeq2/\Delta^2$ is a difference of two numbers close to 1 and carries
an absolute error of order $\eps$. Since $A+B=U_0$, the corresponding error in
$I_c$ is of order $\eps\,S$. It is not amplified, and the $\Lambda$ diagonal
uses $B$, which has no cancellation.

\section{Hermite scheme}

The Hermite solver represents the source over the interval by a cubic B\'ezier
curve in $\tau$ whose end slopes come from the chain-rule (product-rule)
derivatives of $\chi/P_e$ and $S/B_\nu$. The control points are
\begin{equation}
  P_1 = S_u + \tfrac{\Delta}{3}S'_u,\qquad P_2 = S_c - \tfrac{\Delta}{3}S'_c ,
\end{equation}
each clamped into $[\min(S_u,S_c),\max(S_u,S_c)]$. The treatment of slopes,
$\Delta$ (an Euler--Maclaurin trapezoid) and the clamps is unchanged; only the
weights are new. With $u=t/\Delta$ and the Bernstein basis
$b_0=u^3$, $b_1=3u^2(1-u)$, $b_2=3u(1-u)^2$, $b_3=(1-u)^3$,
\begin{equation}
  I_c = E\,I_u + w_0 S_u + w_1 P_1 + w_2 P_2 + w_3 S_c,\qquad
  w_k(\Delta) = \Delta\int_0^1 b_k(u)\,e^{-\Delta u}\,\mathrm{d}u .
  \label{eq:bezier}
\end{equation}

\paragraph{Closed form.}
With $M_n=\int_0^\Delta t^n e^{-t}\mathrm{d}t$, obtained by the recurrence
$M_n = nM_{n-1}-\Delta^nE$ from $M_0=U_0$,
\begin{equation}
  w_0 = \frac{M_3}{\Delta^3},\quad
  w_1 = \frac{3(\Delta M_2-M_3)}{\Delta^3},\quad
  w_2 = \frac{3(\Delta^2M_1-2\Delta M_2+M_3)}{\Delta^3},\quad
  w_3 = M_0-w_0-w_1-w_2 .
\end{equation}
Since $M_n\simeq\Delta^{n+1}/(n+1)$ for small $\Delta$, these combinations lose
all fp32 accuracy below $\Delta\approx2.5$. For $\Delta<0$, where $E>1$ and the
$M_n$ alternate in sign, they lose it everywhere ($w_2$ and $w_3$ are off by
$10^2$ to $10^3$ ulp at $\Delta=-2$).

\paragraph{Series and reflection.}
Expanding the exponential in Eq.~(\ref{eq:bezier}) gives $w_k=\Delta P_k(\Delta)$ with
\begin{equation}
  P_k(\Delta) = \sum_{n\ge0}\frac{(-\Delta)^n}{n!}\int_0^1 b_k(u)\,u^n\,\mathrm{d}u .
\end{equation}
For $\Delta>0$ these are alternating series. $b_0$ and $b_1$ peak at the far
end of the interval ($u=1$ and $u=2/3$), where $e^{-\Delta u}$ is smallest, so
$P_0$ and $P_1$ behave like $e^{-\Delta}$ times a slowly varying factor and their
series cancel strongly. The reflection $u\to1-u$ maps $b_k$ onto $b_{3-k}$:
\begin{equation}
  \frac{w_k(\Delta)}{\Delta} = \int_0^1 b_{3-k}(v)\,e^{-\Delta(1-v)}\mathrm{d}v
  = E\,P_{3-k}(-\Delta) .
\end{equation}
So
\begin{equation}
  w_0 = \Delta E\,P_3(-\Delta),\quad
  w_1 = \Delta E\,P_2(-\Delta),\quad
  w_2 = \Delta\,P_2(\Delta),\quad
  w_3 = \Delta\,P_3(\Delta),
\end{equation}
and for $\Delta>0$ the series for $w_0$ and $w_1$ have terms of one sign. Only
two polynomials are needed: $P_2$ of degree 10 and $P_3$ of degree 9, both
minimax on $|\Delta|\le3$ (each is evaluated at $+\Delta$ and $-\Delta$). The
series are used for $\Delta\le3$ and the closed form above, where the closed
form sums the weights to $U_0$ exactly. Every weight is within 8 ulp of the
exact value for $-3\le\Delta\le87$ (6.5 ulp with fused multiply--add). The fp64
code it replaces used a four-term Taylor series below $\Delta=0.05$, whose
truncation error there is about $10^{-7}$ relative, the same order.

In \fort{bezier3_weights}, with \fort{b3_p2} and \fort{b3_p3} the polynomials
$P_2$ and $P_3$, the closed-form arm is evaluated at \fort{tb = max(t, 3)}, which
equals $\Delta$ wherever that arm is selected and keeps it finite elsewhere:
\begin{minted}{Fortran}
call lw_exp(t, E, M0)
tb  = max(t, 3.0)
tb2 = tb*tb
tb3 = tb2*tb
r3  = 1.0/tb3
M1 = M0 - tb*E
M2 = 2.0*M1 - tb2*E
M3 = 3.0*M2 - tb3*E
b0 = M3*r3
b1 = 3.0*(M2*tb - M3)*r3
b2 = 3.0*(M1*tb2 - 2.0*M2*tb + M3)*r3
ser = t <= 3.0
tE  = t*E
w0 = lw_select(tE*b3_p3(-t), b0, ser)
w1 = lw_select(tE*b3_p2(-t), b1, ser)
w2 = lw_select(t*b3_p2(t), b2, ser)
w3 = lw_select(t*b3_p3(t), M0 - b0 - b1 - b2, ser)
\end{minted}
The interval step is then
\begin{minted}{Fortran}
call bezier3_weights(dt, ex, w0, w1, w2, w3)
Inu = Inuup*ex + w0*ssup + w1*P1 + w2*P2 + w3*ssc
\end{minted}

\section{Ranges and limits}

\begin{itemize}
\item \emph{Maser intervals.} Where $\chi<0$, $\Delta<0$ is capped at
  $-\Delta_{\max}$, set by \fort{maser_dtau_max} (default 2). The series ranges cover this default
  ($|\Delta|\le1$ and $\le2$ for the linear and quadratic schemes, together with
  their closed forms, $\Delta\ge-3$ for the B\'ezier weights). A larger cap
  leaves the linear and quadratic weights accurate, but the B\'ezier weights
  lose accuracy below $\Delta=-3$.
\item \emph{Small intervals.} $|\Delta|$ is held at or above $10^{-20}$, so
  $1/\Delta$ and $\Delta^2$ stay finite and normal.
\item \emph{Thick intervals.} Above $\Delta=87$, $E$ is held at $e^{-87}$; the
  weights are otherwise unaffected.
\end{itemize}

\section{Scope and verification}

\fort{long_char_linear}, \fort{short_char}, \fort{long_char_quadratic} and
\fort{long_char_hermite} now evaluate everything per point in fp32.
\fort{long_char_hermite2} and \fort{short_char_hermite} use the fp32
quadratic moments but still form their weights $A$, $B$, $C$ and, in
\fort{short_char_hermite}, the escape factor $1-\Psi$ in fp64.

The fp32 schemes were compared with the previous fp64 weights on the solar
MARCS model, for the LTE spectrum of the \Code{sed\_4000-9000} benchmark and the
\Code{mg86g} NLTE case, with and without slab decomposition of the atmosphere.
Spectra differ by at most $1.6\times10^{-6}$ relative (median $\sim10^{-7}$),
comparable to the run-to-run spread of threaded NLTE runs. On
\Code{sed\_4000-9000} (one thread) the quadratic solver takes 0.30\,s instead
of 0.88\,s and the Hermite solver 0.43\,s instead of 1.40\,s, because the
frequency loops now vectorise in fp32.

\begin{table}[h]
\centering
\caption{Largest error of the fp32 weights against a quad-precision reference.}
\begin{tabular}{@{}llll@{}}
\toprule
Scheme & Quantities & Range of $\Delta$ & Max.\ error \\
\midrule
shared    & $E$, $U_0$                & $-2\ldots87$ & 2.4 ulp \\
linear    & $E$, $W$, $\Psi$, $1-\Psi$ & $-2\ldots90$ & 4 ulp \\
quadratic & $U_1$, $q$                & $-2\ldots87$ & 3.5 ulp \\
Hermite   & $w_0\ldots w_3$           & $-3\ldots87$ & 7.2 ulp \\
\bottomrule
\end{tabular}
\end{table}

\end{document}
